\documentclass[../../main.tex]{subfiles}
\begin{document}

{\bf [解]}
\paragraph{(1)}

題意の和を$A$とおくと
\begin{align}
 A 
  &=\sum_{n=0}^{99}3^n \\
  &=\dfrac{3^{100}-1}{3-1} \\
  &=\dfrac{1}{2}(3^{100}-1) \\
  &=\dfrac{3^{100}}{2}-\dfrac{1}{2} \label{eq:1}
\end{align}
と変形できる．ここで$A$は整数であり，かつ$3^{100}/2=10^{\ell}, (\ell\in\mathbb{Z})$の形でかけないことから$A$の桁数と
\begin{align}
 B=\dfrac{3^{100}}{2} \label{eq:2}
\end{align}
の桁数は一致する．そこで以下$B$の桁数を考えよう．これを$m$とおくと
\begin{align}
 10^{m-1}\le B<10^m \label{eq:3}
\end{align}
であり，各辺正だから常用対数をとって
\begin{align}
m-1\le\log_{10}B<m,\qquad\log_{10}B=100\log_{10}3-\log_{10}2=47.71-\log_{10}2 \label{eq:4}
\end{align}
を得る．

$\log_{10}x$は単調増加関数であり$0=\log_{10}1<\log_{10}2<\log_{10}3=0.4771$だから
\begin{align}
47.71-0.4771<47.71-\log_{10}2<47.71-0,\quad\text{すなわち}\quad47.2329<\log_{10}B<47.71
\end{align}
となり，これは常に$47<\log_{10}B<48$の範囲にあるから，\cref{eq:4}をみたす$m$は$m=48$のみである．
以上から$B$すなわち$A$は$48$桁である．

\paragraph{(2)}

自然数$n$が与えられた時，自然数$m$を用いて
\begin{align}
 m^2\le n<(m+1)^2 \label{eq:5}
\end{align}
と必ず表せる．この時各辺の平方根をとって
\begin{align}
 m\le\sqrt{n}<m+1
\end{align}
だから，ガウス記号の定義から
\begin{align}
 [\sqrt{n}]=m \label{eq:6}
\end{align}
である．

従って題意の$[\sqrt{n}]$が$n$の約数となるのは，$m$が$n$の約数となる時であるからこれを考えよう．
ここで2以上の自然数$m$に対して
\begin{align}   
   m(m+2) < (m+1)^2 < m(m+3) 
\end{align}
だから，\cref{eq:5}の範囲で$m$が$n$の約数となるのは$m$を固定して$n$を動かして考えると
\begin{align}
 \ n=m^2,\ m(m+1),\ m(m+2)\ \text{の}3\text{通り}
\end{align}
である．
$m=1$の時は常に$m$は$n$の約数で，\cref{eq:5}よりこの時は$1 \le n < 4$となり，$n=1,2,3$の$3$通りである．

考える$n$の最大値は$10000=100^2$だから，$n=10000$も単体で条件を満たす．

以上から$m=1,\dots,99$については各々$3$通りずつ，$m=100$については$n=10000$の$1$通りのみ（範囲が$n\le10000$で打ち切られるため）だから
求める個数は
\begin{align}
 \sum_{m=1}^{99}3+1=297+1=298
\end{align}
である．


{\bf [別解]}

\paragraph{(1)}

$0.4771<0.71<0.4771\times2=\log_{10}9$の各辺に$47$を加えて整理すると
\begin{align}
 \log_{10}10^{47}+\log_{10}3<47.71<\log_{10}10^{47}+\log_{10}9
\end{align}
である．ここで$\log_{10} 3^{100}=47.71$だから
\begin{align}
 \log_{10}10^{47}+\log_{10}3<\log_{10} 3^{100}<\log_{10}10^{47}+\log_{10}9
\end{align}
である．$\log_{10}x$は単調増加だから
\begin{align}
 3\cdot10^{47}<3^{100}<9\cdot10^{47}
\end{align}
を得る．従って$B$すなわち$A$は$48$桁である．

\newpage

\end{document}
